\documentclass[11pt]{article}

\usepackage[T1]{fontenc}
\usepackage[utf8]{inputenc}
\usepackage[margin=1in]{geometry}
\usepackage{booktabs}
\usepackage{amsmath}
\usepackage{amssymb}
\usepackage{xcolor}
\usepackage{enumitem}
\usepackage[colorlinks=true,allcolors=blue!55!black]{hyperref}

\newcommand{\flag}[1]{\texttt{-\/-#1}}
\newcommand{\param}[1]{\subsection*{\texttt{-\/-#1}}}

\title{GAMECA Stage 11 Standout-Analysis Batch Parameters \\[2pt]
\large Reference for the ``Submit batch job'' configuration flags (v0.4.30)}
\author{}
\date{}

\begin{document}
\maketitle

\section*{Overview}
Stage 11 (``standout analysis'') runs a battery of downstream modules on each clustered
TE family — phylogenetics/age estimation, gRNA off-target design, structural variant
detection (transduction, LTR structure), regulatory-context overlays (epigenetics,
CTCF/TAD), cross-assembly and cross-species comparison, structure prediction, and
divergence-landscape / subfamily resolution. Each module exposes its own tunable
parameters; previously these were hard-coded inside \texttt{query.py}'s module table.
As of v0.4.30 the 18 parameters below are surfaced as top-level \texttt{query.py} CLI
flags and as fields in the desktop app's \emph{Submit batch job} panel (grouped under
``Stage 11 — standout analysis modules''). Leaving a field blank uses the underlying
module's own default, so existing behaviour is unchanged unless you opt in.

\bigskip
\noindent\textbf{Conventions used below.}
\emph{Type/options} lists the literal values accepted (free text, integer/float range,
or an enumerated choice list reproduced from the underlying module's \texttt{argparse}
definition). \emph{What it controls} explains which module consumes the value and how.
\emph{Effect of changing it} explains, in plain terms, what differs in the resulting
report when you pick one value over another.

% =====================================================================
\section{Cross-assembly and cross-species comparison}
% =====================================================================

\param{target-assemblies}
\textbf{Module:} \texttt{run\_multiassembly\_liftover.py}\\
\textbf{Type/options:} space-separated list of UCSC assembly codes, e.g.
\texttt{t2t hg19 mm39 panTro6}. Empty $=$ module reports zero target assemblies and
performs no liftover.\\
\textbf{What it controls:} for every code listed, the module fetches/uses the
appropriate \texttt{liftOver} chain file from the source assembly (normally
\texttt{hg38}, taken from the pipeline's \flag{assembly} setting) to the target, lifts
each TE-cluster locus over, and reports which copies have a recoverable orthologous
coordinate in that target genome.\\
\textbf{Effect of changing it:}
\begin{itemize}[leftmargin=1.4em,topsep=2pt,itemsep=2pt]
  \item \texttt{t2t} compares against the telomere-to-telomere CHM13 assembly — the
    one place to look for TE copies that live in regions \emph{absent or collapsed} in
    GRCh38 (centromeres, acrocentric short arms, segmental duplications). This is the
    knob that answers ``does T2T-CHM13 reveal additional/different copies of this
    family?''
  \item \texttt{hg19} compares against the older human reference — useful for tracking
    how a family's annotation/coordinates shift between reference versions, or for
    cross-referencing older literature/datasets that used hg19 coordinates.
  \item Listing a non-human code (\texttt{mm39}, \texttt{panTro6}, \ldots) repurposes
    the \emph{same} liftover machinery as a coarse cross-species presence/absence check
    (a lighter-weight companion to the dedicated ortholog module below).
  \item Listing more assemblies simply runs more liftover passes and adds more columns
    to the comparison report — there is no quality trade-off, only runtime.
\end{itemize}

\param{ortholog-species}
\textbf{Module:} \texttt{run\_ortholog\_insertion.py}\\
\textbf{Type/options:} space-separated list of UCSC assembly codes for
\emph{non-human} species, e.g. \texttt{panTro6} (chimpanzee), \texttt{rheMac10}
(rhesus macaque), \texttt{mm39} (mouse). Empty $=$ ``0 species'' and the module is
effectively a no-op.\\
\textbf{What it controls:} for each species, the module auto-downloads the
\texttt{hg38}$\to$species chain file from UCSC and tests whether each human TE
insertion has an orthologous (syntenic) locus in that species' genome.\\
\textbf{Effect of changing it:} this is the dating tool for individual insertions —
\begin{itemize}[leftmargin=1.4em,topsep=2pt,itemsep=2pt]
  \item If an insertion \emph{is} present at the orthologous locus in chimpanzee
    (\texttt{panTro6}), it predates the human/chimp split ($\sim$6--7 Mya); if it is
    present in macaque (\texttt{rheMac10}) it predates the Old-World-monkey split
    ($\sim$25--30 Mya); presence in mouse (\texttt{mm39}) implies an ancient,
    pre-Boreoeutherian element.
  \item Conversely, an insertion \emph{absent} from the orthologous locus in a species
    but present in human is \emph{lineage-specific} — i.e., it inserted after that
    species diverged from the human lineage. Choosing more (and more divergent) species
    lets the module triangulate a tighter age bracket and flag human-specific
    insertions.
  \item Species closer to human (chimp, macaque) are faster to lift over and have
    higher-quality chains; very distant species (mouse and beyond) take longer and
    yield lower mapping rates simply because so much intervening sequence has
    diverged/rearranged.
\end{itemize}

\param{liftover-cmd}
\textbf{Modules:} \texttt{run\_ortholog\_insertion.py},
\texttt{run\_multiassembly\_liftover.py}\\
\textbf{Type/options:} free-text path or command name. Default: \texttt{liftOver}
(resolved via \texttt{PATH}).\\
\textbf{What it controls:} which UCSC \texttt{liftOver} executable both
liftover-dependent modules invoke.\\
\textbf{Effect of changing it:} purely an environment knob — set it to an absolute
path (e.g. \texttt{/opt/tools/liftOver}) when the binary is installed somewhere not on
the remote host's \texttt{PATH}. It does not change any analysis result, only whether
the liftover steps can run at all.

% =====================================================================
\section{Regulatory and chromatin-context overlays}
% =====================================================================

\param{epigenetic-preset}
\textbf{Module:} \texttt{run\_epigenetic\_overlay.py}\\
\textbf{Type/options (choice):} \texttt{K562}, \texttt{GM12878}, \texttt{HeLa-S3},
\texttt{IMR-90}, \texttt{H1-hESC}. Blank $=$ ``0 tracks used'' (overlay skipped).\\
\textbf{What it controls:} which ENCODE reference cell line's curated panel of
regulatory tracks (DNase hypersensitivity, histone marks such as H3K27ac/H3K4me1/
H3K4me3, chromatin-state segmentations, \ldots) the module intersects TE loci against.\\
\textbf{Effect of changing it — what each cell line represents:}
\begin{itemize}[leftmargin=1.4em,topsep=2pt,itemsep=2pt]
  \item \textbf{K562}: chronic myelogenous leukaemia line — erythroid/myeloid lineage,
    very widely profiled by ENCODE (deepest, most complete track panel).
  \item \textbf{GM12878}: B-lymphoblastoid line — immune/lymphoid lineage; the
    canonical line for 3D-genome (Hi-C) and CTCF studies.
  \item \textbf{HeLa-S3}: cervical-adenocarcinoma line — generic transformed epithelial
    chromatin landscape.
  \item \textbf{IMR-90}: primary lung fibroblast — a non-cancer, mesenchymal reference,
    useful as a contrast to the cancer lines above.
  \item \textbf{H1-hESC}: human embryonic stem cell line — the developmental/pluripotent
    reference; the only choice that reflects an undifferentiated chromatin state.
\end{itemize}
Because enhancers, promoters and open chromatin are substantially cell-type specific,
the preset you choose determines \emph{which} TE copies show up as ``sitting in active
regulatory sequence'' — a TE can be silent in K562 chromatin and active in H1-hESC, or
vice versa. Pick the line that best matches the biological context you care about (e.g.
H1-hESC if you are studying TE co-option in pluripotency/early development; GM12878 or
K562 for blood/immune contexts).

\param{ctcf-preset}
\textbf{Module:} \texttt{run\_ctcf\_tad.py}\\
\textbf{Type/options (choice):} \texttt{K562}, \texttt{GM12878}, \texttt{HeLa-S3}.
Blank $=$ no CTCF peak set loaded.\\
\textbf{What it controls:} which ENCODE CTCF ChIP-seq peak catalogue TE loci are
checked against, in addition to the module's intrinsic CTCF-motif scan.\\
\textbf{Effect of changing it:} CTCF is the principal architectural/insulator protein
that defines genome-folding boundaries; its binding is partly cell-type specific. The
preset determines whether a TE-borne CTCF motif is reported merely as a
\emph{sequence-level candidate} or as an \emph{experimentally observed binding site} in
that line — GM12878 is the deepest-profiled line for CTCF/3D-genome work and is a
reasonable default choice if you have no other preference; K562 and HeLa-S3 give a
cancer-lineage comparison point.

\param{tads-preset}
\textbf{Module:} \texttt{run\_ctcf\_tad.py}\\
\textbf{Type/options (choice):} \texttt{K562}, \texttt{GM12878}, \texttt{IMR90}.
Blank $=$ no TAD boundary map loaded.\\
\textbf{What it controls:} which Hi-C–derived topologically-associating-domain (TAD)
boundary map is used to test whether TE loci sit at or near domain boundaries (where
TE-borne CTCF/insulator activity could plausibly contribute to 3D-genome architecture,
or where a TE insertion could disrupt an existing boundary).\\
\textbf{Effect of changing it:} TAD structure, like CTCF occupancy, differs somewhat
between cell types/lineages. \texttt{GM12878} has the best-characterised Hi-C maps;
\texttt{K562} offers a myeloid/cancer comparison; \texttt{IMR90} (note: no hyphen, unlike
the epigenetic-overlay \texttt{IMR-90}) is the only primary, non-cancer fibroblast option
and is informative if you want a ``normal tissue'' 3D-genome reference rather than a
transformed line.

% =====================================================================
\section{CRISPR gRNA design and off-target scoring}
% =====================================================================

\param{grna-cas}
\textbf{Module:} \texttt{run\_grna\_offtarget.py}\\
\textbf{Type/options (choice), with PAM motif:}
\begin{center}
\begin{tabular}{@{}lll@{}}
\toprule
Value & PAM & Notes \\
\midrule
\texttt{SpCas9} (default) & \texttt{NGG} & Standard \emph{S. pyogenes} Cas9; most common, best characterised \\
\texttt{SaCas9} & \texttt{NNGRRT} & Smaller \emph{S. aureus} Cas9; fits in AAV vectors for delivery \\
\texttt{Cas12a} & \texttt{TTTV} & (Cpf1) staggered-cut, T-rich PAM; suits AT-rich targets \\
\texttt{SpRY} & \texttt{NRN} & Near-PAMless engineered SpCas9 variant; broadest targeting range \\
\bottomrule
\end{tabular}
\end{center}
\textbf{What it controls:} which PAM motif the module scans the TE consensus/copies for
when enumerating candidate guide sites, and therefore the entire downstream candidate
pool and off-target search.\\
\textbf{Effect of changing it:} the PAM requirement is the single biggest determinant
of \emph{how many} candidate guides exist and \emph{where} they can sit. \texttt{NGG}
(SpCas9) occurs roughly every $\sim$8--16\,bp in human sequence and is the
``workhorse'' choice. \texttt{NNGRRT} (SaCas9) is rarer and yields fewer candidates but
matches a smaller, AAV-deliverable enzyme used in therapeutic contexts. \texttt{TTTV}
(Cas12a) is biased toward AT-rich stretches — useful if the TE consensus is GC-poor —
and produces staggered double-strand breaks rather than blunt ends. \texttt{NRN}
(SpRY) is by far the most permissive PAM and will surface the largest candidate pool
and the densest possible tiling of guides across a locus, at the cost of also
broadening the genomic off-target search space (more places in the genome satisfy a
near-PAMless requirement, so off-target lists tend to be longer).

\param{grna-max-mm}
\textbf{Module:} \texttt{run\_grna\_offtarget.py}\\
\textbf{Type/options (choice):} \texttt{0}, \texttt{1}, \texttt{2}. Default:
\texttt{2}.\\
\textbf{What it controls:} the maximum number of base mismatches tolerated when the
module searches the genome for sequences that a candidate guide could also bind
(potential off-target sites).\\
\textbf{Effect of changing it:}
\begin{itemize}[leftmargin=1.4em,topsep=2pt,itemsep=2pt]
  \item \texttt{0} — only exact 20-mer matches elsewhere in the genome are counted as
    off-targets. Fastest to compute, but under-estimates risk: real Cas9 tolerates
    several mismatches (especially distal from the PAM), so this setting will miss
    most genuine off-target liabilities.
  \item \texttt{1} — also counts near-exact matches differing by a single base; a
    middle ground.
  \item \texttt{2} (default) — also counts sites differing by up to two bases, which is
    closer to the empirically observed mismatch tolerance of SpCas9 \emph{in vivo} and
    gives the most realistic (if more conservative and slower-to-compute) off-target
    risk picture, especially important when many copies of a repetitive TE family share
    near-identical sequence across the genome.
\end{itemize}

\param{grna-background}
\textbf{Module:} \texttt{run\_grna\_offtarget.py}\\
\textbf{Type/options:} free-text path to a FASTA file. Default: empty (no background
genome used).\\
\textbf{What it controls:} when supplied, every candidate guide is \emph{additionally}
scored for off-target sites in this second genome/sequence set, on top of the primary
reference.\\
\textbf{Effect of changing it — this is the ``allele-aware'' knob:} point it at, e.g.,
a second haplotype/allele assembly, a different mouse strain, or a related species'
genome, and the report will flag guides that look clean against the primary reference
but would still cut in the background sequence (or, just as usefully, guides that are
specific to one allele/strain and \emph{not} the other — which can be exploited for
allele-selective editing). Leaving it blank simply skips this second pass.

% =====================================================================
\section{Structure prediction}
% =====================================================================

\param{colabfold-cmd}
\textbf{Module:} \texttt{run\_fold\_prediction.py}\\
\textbf{Type/options:} free-text path or command name for the \texttt{colabfold\_batch}
executable. Default: empty $\Rightarrow$ the module reports
\texttt{colabfold\_batch not found} and skips structure prediction entirely.\\
\textbf{What it controls:} whether/how the module invokes ColabFold (a local
AlphaFold2 front-end) to predict the 3-D structure of intact open reading frames found
in TE consensus sequences (e.g. reverse transcriptase, integrase/endonuclease,
gag/pol/env-like domains).\\
\textbf{Effect of changing it:} this is purely an enable/path switch — there is no
``analysis'' difference between values, only whether the feature runs at all. Set it to
the absolute path of your \texttt{colabfold\_batch} installation (e.g.
\texttt{/opt/colabfold/colabfold\_batch}) to turn structure prediction \emph{on} and
get per-cluster predicted folds (with confidence/pLDDT plots) added to the report;
leave it blank to skip the (computationally heavy) step entirely.

% =====================================================================
\section{Phylogenetics, divergence-based age, and master-element identification}
% =====================================================================

\param{subst-rate}
\textbf{Module:} \texttt{run\_phylo\_analysis.py}\\
\textbf{Type/options:} positive float, substitutions per site per year. Default:
\texttt{2.2e-9} (the commonly cited human neutral substitution rate).\\
\textbf{What it controls:} the molecular-clock conversion factor used to translate
observed sequence divergence (branch lengths on the phylogenetic tree / \% divergence
from consensus) into an estimated insertion age in years.\\
\textbf{Effect of changing it:} every age estimate the module reports is
(to first order) inversely proportional to this number — doubling the rate halves all
inferred ages, and vice versa. \textbf{The value should match the species/lineage being
analysed}: the human rate ($\sim$2.2$\times10^{-9}$) is not appropriate for, say, a
mouse genome (rodents have a substantially higher substitution rate, commonly cited
around $4$--$5\times10^{-9}$); using the wrong rate systematically biases every
downstream age/timeline statement in the report. Override this whenever
\flag{assembly} points at a non-human genome.

\param{clock-divisor}
\textbf{Module:} \texttt{run\_phylo\_analysis.py}\\
\textbf{Type/options:} positive float (a unitless correction factor). Default:
\texttt{2}.\\
\textbf{What it controls:} a secondary scaling factor applied on top of
\flag{subst-rate} — e.g. to fold in a generation-time correction or a
lineage-specific adjustment to the effective clock rate.\\
\textbf{Effect of changing it:} \emph{larger} divisor $\Rightarrow$ the same observed
divergence is interpreted as accumulating \emph{faster} per calendar year, so the
inferred insertion age comes out \emph{younger}; \emph{smaller} divisor (down to
\texttt{1}, i.e. no correction) $\Rightarrow$ the same divergence implies an
\emph{older} age. Use this to apply published lineage-specific clock corrections without
having to recompute \flag{subst-rate} from scratch.

\param{intact-orf-aa}
\textbf{Module:} \texttt{run\_phylo\_analysis.py}\\
\textbf{Type/options:} positive integer, amino-acid length. Default: \texttt{100}.\\
\textbf{What it controls:} the minimum length an open reading frame must reach, within
a TE copy's sequence, to be classified as ``intact'' — i.e. the threshold used to flag
a copy as a candidate \emph{master/source/active element} (one that plausibly retains
enough protein-coding capacity to still be capable of retrotransposition).\\
\textbf{Effect of changing it:}
\begin{itemize}[leftmargin=1.4em,topsep=2pt,itemsep=2pt]
  \item \emph{Lower} the threshold (e.g. to 50\,aa) $\Rightarrow$ more copies pass the
    bar and get reported as ``intact/candidate master elements'' — more permissive,
    higher sensitivity, but also a higher chance of counting truncated, non-functional
    fragments as candidates (false positives).
  \item \emph{Raise} it (e.g. to 300\,aa, closer to a full functional domain such as
    reverse transcriptase) $\Rightarrow$ stricter, fewer candidates reported, but those
    that remain are far more likely to be genuinely capable of encoding a functional
    protein — i.e. higher specificity for true master/source elements.
\end{itemize}

\param{mafft-cmd}
\textbf{Module:} \texttt{run\_phylo\_analysis.py}\\
\textbf{Type/options:} free-text path or command name. Default: \texttt{mafft}
(resolved via \texttt{PATH}).\\
\textbf{What it controls:} which \texttt{mafft} multiple-sequence-alignment binary the
module calls to build the alignments that the phylogenetic tree and all
divergence/age statistics are derived from.\\
\textbf{Effect of changing it:} an environment knob, not an analysis knob — point it at
an absolute path (e.g. \texttt{/opt/homebrew/bin/mafft}) when \texttt{mafft} is not on
the remote host's \texttt{PATH}. (This directly addresses the ``mafft not found''
fallback warning seen in Stage 11 logs — supplying a working path here lets the
alignment-dependent phylogenetics actually run instead of falling back.)

% =====================================================================
\section{LTR structural annotation}
% =====================================================================

\param{min-ltr-identity}
\textbf{Module:} \texttt{run\_ltr\_struct.py}\\
\textbf{Type/options:} float in $[0,1]$ (a sequence-identity fraction). Default:
\texttt{0.65}.\\
\textbf{What it controls:} the minimum sequence identity required between a TE copy's
5$'$ and 3$'$ long terminal repeats (LTRs) for that copy to be classified as having
intact/well-matched LTR structure.\\
\textbf{Effect of changing it:} when an LTR retrotransposon first integrates, its two
flanking LTRs are exact copies of one another (both derived from the same template
during reverse transcription); they then diverge independently over time through
ordinary mutation accumulation. \textbf{LTR–LTR identity is therefore itself a proxy
for insertion age} — high identity means young/recent, low identity means old/decayed.
\begin{itemize}[leftmargin=1.4em,topsep=2pt,itemsep=2pt]
  \item \emph{Raise} the threshold (e.g. to \texttt{0.95}) $\Rightarrow$ only the
    youngest, most pristine copies are called ``structurally intact'' — a stricter,
    higher-confidence (but smaller) set, useful for identifying very recent/active
    insertions.
  \item \emph{Lower} it (e.g. to \texttt{0.50}) $\Rightarrow$ more, older, more-degraded
    copies are still recognised as having LTR structure — broader coverage of the
    family's structural history, at the cost of including more decayed/ambiguous calls.
\end{itemize}

% =====================================================================
\section{3$'$ transduction detection}
% =====================================================================

\param{tail-bp}
\textbf{Module:} \texttt{run\_transduction.py}\\
\textbf{Type/options:} positive integer, base pairs. Default: \texttt{150}.\\
\textbf{What it controls:} the size of the genomic window immediately downstream
(3$'$) of each TE copy that the module scans for ``transduced'' sequence — host-genome
sequence that gets carried along into a new insertion site because the
retrotransposition machinery read through the element's polyadenylation signal into
adjacent unique sequence (a well-documented hallmark of active L1/Alu-type elements,
``3$'$ transduction'').\\
\textbf{Effect of changing it:}
\begin{itemize}[leftmargin=1.4em,topsep=2pt,itemsep=2pt]
  \item \emph{Larger} window (e.g. 300--500\,bp) $\Rightarrow$ the module can detect
    longer transduced segments, but a wider scan window also raises the chance of
    matching unrelated, coincidentally-shared flanking sequence — more false positives.
  \item \emph{Smaller} window (e.g. 50--100\,bp) $\Rightarrow$ a more conservative,
    higher-specificity scan that will still catch short transductions but may miss
    longer ones.
\end{itemize}
The default of 150\,bp is a reasonable middle ground for human L1-class transductions,
which are typically in the tens-to-low-hundreds of base pairs.

% =====================================================================
\section{Antisense / bidirectional promoter activity}
% =====================================================================

\param{promoter-bp}
\textbf{Module:} \texttt{run\_antisense\_promoter.py}\\
\textbf{Type/options:} positive integer, base pairs. Default: \texttt{200}.\\
\textbf{What it controls:} the size of the window, anchored on each TE copy, that the
module scans — in \emph{both} sense and antisense orientation — for promoter-like
sequence signatures (testing whether the element's internal sequence can drive
transcription in the reverse direction; a documented mechanism by which TEs spawn novel
transcripts and rewire nearby gene-expression programmes, i.e. ``antisense/bidirectional
promoter activity'').\\
\textbf{Effect of changing it:}
\begin{itemize}[leftmargin=1.4em,topsep=2pt,itemsep=2pt]
  \item \emph{Larger} window $\Rightarrow$ more sequence context examined per copy —
    can pick up promoter-like motifs that sit slightly outside the canonical TE
    boundary (e.g. in adjacent genomic sequence shaped by the insertion), but also
    increases the chance of flagging motifs that have nothing to do with the TE itself.
  \item \emph{Smaller} window $\Rightarrow$ a tighter scan focused on signal intrinsic
    to the element's own sequence — more specific, but may under-count
    promoter activity that depends on the local genomic context just outside the TE.
\end{itemize}

% =====================================================================
\section{Divergence landscape and subfamily resolution}
% =====================================================================

\param{cpg-omega}
\textbf{Modules:} \texttt{run\_divergence.py}, \texttt{run\_subfamily.py}\\
\textbf{Type/options:} positive float (a unitless rate-multiplier). Default:
\texttt{10.0}.\\
\textbf{What it controls:} the correction factor applied, in the Kimura
two-parameter divergence calculation, to up-weight the expected mutation rate at CpG
dinucleotides relative to other sites. (Methylated cytosines at CpG sites spontaneously
deaminate to thymine at a substantially elevated rate — a textbook molecular-evolution
correction; the literature default multiplier is commonly taken as $\sim$10$\times$ the
background rate.)\\
\textbf{Effect of changing it:}
\begin{itemize}[leftmargin=1.4em,topsep=2pt,itemsep=2pt]
  \item In \texttt{run\_divergence.py}, this factor feeds directly into the
    ``repeat-landscape'' Kimura-divergence histogram — the classic plot of how many TE
    copies fall into each \%-divergence-from-consensus bin (a proxy for the family's
    age distribution / activity history). A higher \texttt{cpg-omega} more aggressively
    discounts CpG-site differences when computing each copy's corrected divergence,
    which can shift and sharpen the inferred age distribution (especially for
    CpG-rich consensus sequences); a lower value (down to \texttt{1.0}, i.e. no CpG
    correction) treats all sites uniformly and tends to inflate apparent divergence/age
    for CpG-rich families.
  \item In \texttt{run\_subfamily.py}, the same corrected-divergence metric is one of
    the inputs to automatic subfamily clustering — changing \texttt{cpg-omega} can
    therefore shift which copies are judged ``close enough'' to group into the same
    putative subfamily, and in turn how many subfamilies the module resolves.
\end{itemize}
Use the default (\texttt{10.0}) unless you have a specific reason (e.g. a published,
lineage-specific CpG-mutation-rate estimate) to override it — both modules' outputs are
designed around that standard correction.

\end{document}
